\section*{Capítulo \ref{ch:b3:product} --- Medidas produto, Fubini, mudança de variáveis}

\begin{solution}{exo:b3:product:1}
$\Delta $ está fechado em $\intcc01^2$, daí Borel, e $\mathcal
B(\intcc01^2)$ é o produto $\sigma $-álgebra
(\cref{prop:b3:product:sections}(b)). Iterando de uma maneira:
\[
\int_{\intcc01}\Bigl(\int \mathbf 1_\Delta(x,y)
\,\dd\lambda(y)\Bigr)\dd\mu(x)
= \int \lambda(\{x\})\,\dd\mu(x) = 0 ;
\]
do outro jeito:
\[
\int_{\intcc01}\Bigl(\int\mathbf
1_\Delta(x,y)\,\dd\mu(x)\Bigr)\dd\lambda(y)
= \int \mu(\{y\})\,\dd\lambda(y) = \int 1\,\dd\lambda = 1 .
\]
A hipótese falha é $\sigma $-finitude da contagem
medida $\mu $ no incontável $\intcc01$: não contável
família de finitos $\mu $ conjuntos cobre isso.
\end{solution}

\begin{solution}{exo:b3:product:2}
Sobre $\intcc0A\times\intoo0{+\infty}$:
$\int_0^A\int_0^\infty\eu^{-xy}\abs{\sin x}\,\dd y\,\dd x =
\int_0^A\frac{\abs{\sin x}}x\dd x \leq A < \infty $ (Tonelli para
o valor absoluto): Fubini aplica-se, e uma vez que
$\int_0^\infty\eu^{-xy}\dd y = \frac1x $,
\[
\int_0^A\frac{\sin x}x\dd x
= \int_0^\infty\Bigl(\int_0^A\eu^{-xy}\sin x\,\dd
x\Bigr)\dd y
= \int_0^\infty \frac{1 - \eu^{-Ay}(\cos A + y\sin A)}{1 +
y^2}\,\dd y
\]
(a integral interna: $\operatorname{Im}\int_0^A\eu^{(\iu -
y)x}\dd x $, calculado diretamente). Como $ A \to \infty $, o
prazo de correção é limitado por $\int_0^\infty\eu^{-Ay}\frac{1 +
y}{1 + y^2}\dd y \leq \frac32\int_0^\infty\eu^{-Ay}\dd y =
\frac3{2A} \to 0$; o termo principal é
$\int_0^\infty\frac{\dd y}{1+y^2} = \frac\pi2$. Por isso
$\int_0^\infty\frac{\sin x}x\dd x = \frac\pi2$ --- Dirichlet's
integral de Fubini.
\end{solution}

\begin{solution}{exo:b3:product:3}
(a) Bolo em camadas (\cref{prop:b3:product:layercake}): $\int
f\,\dd\mu = \int_0^\infty\mu(f > t)\,\dd t $, e $ t \mapsto
\mu(f > t)$ é não crescente. Sobre $[n-1, n]$: $\mu(f \geq n)
\leq \mu(f > t) \leq \mu(f > n - 1) \leq \mu(f \geq n - 1)$;
somando as integrais nos intervalos unitários:
\[
\sum_{n\geq1}\mu(f \geq n) \leq \int f\,\dd\mu \leq
\sum_{n\geq1}\mu(f \geq n - 1) = \mu(f \geq 0) +
\sum_{n\geq1}\mu(f\geq n) \leq \mu(X) + \sum_{n\geq1}\mu(f\geq
n).
\]

(b) Aplicar (a) a $\abs f $: finitude da integral e de
as séries são equivalentes (o extra $\mu(X)$ é finito).
\end{solution}

\begin{solution}{exo:b3:product:4}
Desde $ f(x,y) = \partial_y\bigl(\frac{y}{x^2+y^2}\bigr)$ para
$ x \ne 0$:
\[
\int_0^1 f(x, y)\,\dd y = \frac{1}{x^2 + 1}
\ \Longrightarrow\
\int_0^1\!\!\int_0^1 f\,\dd y\,\dd x =
\int_0^1\frac{\dd x}{1 + x^2} = \frac\pi4 ;
\]
pela antisimetria $ f(y,x) = -f(x,y)$, a outra ordem dá
$-\frac\pi4$. Valores absolutos: para $0 < y < x $,
\[
\int_0^x f(x,y)\,\dd y = \Bigl[\frac{y}{x^2 +
y^2}\Bigr]_0^x = \frac1{2x},
\quad\text{and } f \geq 0 \text{ there, so}\quad
\int_0^1\!\!\int_0^1\abs f \geq \int_0^1\frac{\dd x}{2x} =
+\infty .
\]
Nenhuma contradição com Fubini: sua hipótese $ f \in L^1$
falha, e as duas integrais iteradas são simplesmente duos números.
\end{solution}

\begin{solution}{exo:b3:product:5}
(um) $(\mathbf 1_{\intcc01}*\mathbf 1_{\intcc01})(x) =
\lambda\bigl(\intcc01\cap\intcc{x-1}x\bigr)$: $0$ para $ x
\notin \intcc02$, $ x $ para $0 \leq x \leq 1$, $2 - x $ para $1
\leq x \leq 2$: a tenda. Convolver novamente dá uma $\mathcal
C^1$ colisão quadrática por partes $\intcc03$ (o quadrático
B-spline): cada convolução ganha um \omterm{def:b3:galois:extension}{grau} de suavidade ---
o princípio de suavização por trás \cref{ch:b3:lp}são apaziguadores.

(b) Se $ x \notin \overline{\operatorname{supp}f +
\operatorname{supp}g}$, há uma bola ao redor $ x $ disjunto
da soma definida; para $ y \in \operatorname{supp}g $, $ x - y
\notin\operatorname{supp}f $, então o integrando desaparece
de forma idêntica: $ f * g = 0$ aproximar $ x $.

(c) Para $ x_n \to x $: $(f*g)(x_n) = \int
f(y)g(x_n - y)\,\dd y $; os integrandos convergem ponto a ponto
(continuidade de $ g $) e são dominados por $\norm
g_\infty\,\abs f \in L^1$: DCT dá $(f*g)(x_n) \to (f*g)(x)$.
\end{solution}

\begin{solution}{exo:b3:product:6}
(a) Por Fubini e indução, cortando ao longo do último
coordenada:
\[
\lambda_d(\Delta_d) = \int_0^1
\lambda_{d-1}\bigl((1 - t)\,\Delta_{d-1}\bigr)\,\dd t
= \lambda_{d-1}(\Delta_{d-1})\int_0^1(1 - t)^{d-1}\dd t
= \frac{\lambda_{d-1}(\Delta_{d-1})}{d},
\]
usando a regra de dilatação $\lambda_{d-1}(\rho A) =
\rho^{d-1}\lambda_{d-1}(A)$
(\cref{thm:b3:product:linearchange}); com
$\lambda_1(\Delta_1) = 1$: volume $\frac1{d!}$.

(b) $ v_2 = \pi/\Gamma(2) = \pi $; $ v_3 =
\pi^{3/2}/\Gamma(\frac52) = \pi^{3/2}/(\frac32\cdot\frac12
\sqrt\pi) = \frac{4\pi}3$. O elipsóide é $ T(B(0,1))$ com
$ T = \operatorname{diag}(a_1, \dots, a_d)$:
\cref{thm:b3:product:linearchange} dá volume
$ v_d\prod a_i $.
\end{solution}

\begin{solution}{exo:b3:product:7}
Em $\R^2$, \omterm{ex:b3:product:polar}{coordenadas polares} (\cref{ex:b3:product:polar}):
\[
\int_{B(0,1)}\frac{\dd x\dd y}{(x^2+y^2)^s}
= 2\pi\int_0^1 r^{1 - 2s}\,\dd r,
\qquad
\int_{\R^2\setminus B(0,1)} = 2\pi\int_1^\infty r^{1-2s}\dd r:
\]
finito se $1 - 2s > -1$ ($ s < 1$), resp.\ $1 - 2s < -1$ ($ s >
1$). Em $\R^d $, evite coordenadas esféricas com a camada
bolo: $\lambda_d(\{\norm x^{-2s} > t\}\cap B(0,1)) =
\lambda_d(B(0, \min(1, t^{-1/2s}))) = v_d\min(1, t^{-d/2s})$,
e $\int_0^\infty v_d\min(1, t^{-d/(2s)})\dd t < \infty $ se
$\frac d{2s} > 1$, ou seja\ $ s < \frac d2$; a integral externa
converge se e se $ s > \frac d2$ (mesmo cálculo no
região complementar).
\end{solution}

\begin{solution}{exo:b3:product:8}
Por Tonelli (integrandos positivos) e a mudança de variáveis
$(x, y) = \Phi(u, v) = (uv,\ u(1-v))$, um $\mathcal C^1$
difeomorfismo de $\intoo0\infty\times\intoo01$ a o aberto
quadrante com
\[
\det D\Phi = \det\begin{pmatrix} v & u\\ 1 - v & -u
\end{pmatrix} = -uv - u(1 - v) = -u,
\qquad \abs{\det} = u :
\]
\[
\Gamma(p)\Gamma(q) =
\iint x^{p-1}y^{q-1}\eu^{-x-y}\dd x\,\dd y
= \iint (uv)^{p-1}\bigl(u(1{-}v)\bigr)^{q-1}\eu^{-u}\,u\,
\dd u\,\dd v
= \Gamma(p + q)\,B(p, q).
\]
Substituindo $ t = \sin^2\theta $ em $ B(p,q)$ dá
$2\int_0^{\pi/2}\sin^{2p-1}\theta\cos^{2q-1}\theta\,\dd\theta =
B(p, q)$. Wallis: $ W_n =
\int_0^{\pi/2}\sin^n\theta\,\dd\theta = \frac12B\bigl(\frac{n +
1}2, \frac12\bigr) = \frac{\Gamma(\frac{n+1}2)\sqrt\pi}
{2\,\Gamma(\frac n2 + 1)}$ --- por exemplo\ $ W_{2n} =
\frac\pi2\cdot\frac{(2n)!}{4^n(n!)^2}$ usando
$\Gamma(n + \frac12) = \frac{(2n)!}{4^nn!}\sqrt\pi $.
\end{solution}

\begin{solution}{exo:b3:product:9}
Indicadores: $\int\mathbf 1_B\,\dd(T_*\mu) = T_*\mu(B) =
\mu(T^{-1}B) = \int\mathbf 1_B\circ T\,\dd\mu $; linearidade
estende-se ao \omterm{def:b3:groups:simple}{simples} $ g $, MCT para $ g \geq 0$ --- a transferência
fórmula. É puramente formal: reexpressa integrais contra
$ T_*\mu $ mas não diz nada sobre \emph{o que} $ T_*\mu $ é. O
conteúdo de \cref{thm:b3:product:linearchange} e
\cref{thm:b3:product:changeofvar} é a identificação
\[
\Phi_*\bigl(\lambda_d\restriction_U\bigr) =
\abs{\det D\Phi^{-1}}\,\lambda_d\restriction_V
\quad\text{(a density measure)},
\]
ou seja, um cálculo do pushforward de Medida Lebesgue ---
a entrada analítica é a geometria diferencial de $\Phi $,
não formalismo teórico da medida.
\end{solution}

\begin{solution}{exo:b3:product:10}
Por Tonelli, a fatoração gaussiana, então com $ G_1 =
\int_\R\eu^{-s^2}\dd s = \sqrt\pi $ e $\int_\R
s^2\eu^{-s^2}\dd s = \frac{\sqrt\pi}2$ (integrar por partes):
\[
\int_{\R^d}x_1^2\,\eu^{-\norm x^2}\dd x =
\frac{\sqrt\pi}2\;\pi^{(d-1)/2} = \frac{\pi^{d/2}}2,
\qquad
\int_{\R^d}\norm x^2\eu^{-\norm x^2}\dd x =
d\cdot\frac{\pi^{d/2}}2
\]
(por simetria, $\norm x^2 = \sum_ix_i^2$ contribui $ d $ igual
termos --- a verificação de consistência). Para o normalizado medida
$\pi^{-d/2}\eu^{-\norm x^2}\dd x $, o segundo momento é
$\frac d2$.
\end{solution}

\begin{solution}{exo:b3:product:11}
(um) $ H = \Phi^{-1}(\intoo0\infty)$ para $\Phi(x, y) = f(x) -
y $ cruzado com $\{y > 0\}$: \omterm{def:b3:lebesgue:measurable}{mensurável}, desde $(x, y)
\mapsto f(x)$ e $(x,y)\mapsto y $ são \omterm{def:b3:lebesgue:measurable}{mensurável} no
produto (composições com as projeções). O
$ x $-seção de $ H $ é $\intoo0{f(x)}$, de medida $ f(x)$:
Tonelli integra as seções,
\[
\lambda_{d+1}(H) = \int_{\R^d}\lambda_1\bigl(\intoo0{f(x)}
\bigr)\,\dd x = \int_{\R^d}f\,\dd\lambda_d .
\]

(b) O gráfico é $\{(x,y) : y \geq f(x)\} \cap \{y \leq
f(x)\}$, \omterm{def:b3:lebesgue:measurable}{mensurável}; isso é $ x $-seções são singletons, de
medida $0$: Tonelli dá $\lambda_{d+1}(\text{graph}) =
\int 0 = 0$.

(c) $ S^{d-1}$ é a união dos dois gráficos $ y =
\pm\sqrt{1 - \abs{x'}^2}$ sobre a bola unitária de $\R^{d-1}$
(dividindo a última coordenada): uma união de dois conjuntos nulos por
(b), nulo.
\end{solution}

\begin{solution}{exo:b3:product:12}
Sobre $\intoo01^2$, $\frac1{1 - xy} = \sum_{n\geq0}(xy)^n $ com
termos não negativos: Tonelli permite integração termo a termo,
\[
\iint\frac{\dd x\,\dd y}{1 - xy}
= \sum_{n\geq0}\Bigl(\int_0^1x^n\dd x\Bigr)
\Bigl(\int_0^1y^n\dd y\Bigr)
= \sum_{n\geq0}\frac1{(n+1)^2} = \zeta(2) .
\]
Para o caso alternado, $\frac1{1 + xy} =
\sum_n(-1)^n(xy)^n $ não é uma série positiva; mas o
integral do \emph{absoluto} série é $\zeta(2) <
\infty $, então Fubini (integrabilidade agora estabelecido) aplica-se:
$\iint\frac{\dd x\dd y}{1 + xy} =
\sum_n\frac{(-1)^n}{(n+1)^2}$. A identidade da série:
\[
\sum_{n\geq1}\frac{(-1)^{n-1}}{n^2}
= \sum_{n\geq1}\frac1{n^2} - 2\sum_{k\geq1}\frac1{(2k)^2}
= \zeta(2) - \frac{\zeta(2)}2 = \frac{\zeta(2)}2 .
\]
Com $\zeta(2) = \frac{\pi^2}6$ (\cref{pb:b3:spectral:1}):
as integrais são $\frac{\pi^2}6$ e $\frac{\pi^2}{12}$.
A positividade de Tonelli foi o jogo inteiro no primeiro
cálculo --- não integrabilidade verifique necessário antes
trocando; no segundo, positividade do absoluto
série é o que \emph{certifica} integrabilidade então isso
Fubini poderá concorrer no assinado.
\end{solution}

\begin{solution}{pb:b3:product:1}
\textbf{1.} Com $ x = t + \sqrt t\,u $ ($\dd x = \sqrt
t\,\dd u $; $ x $ varia acima $\intoo0\infty $ como $ u $ varia acima
$\intoo{-\sqrt t}\infty $):
\[
\Gamma(t{+}1) = \int_0^\infty x^t\eu^{-x}\dd x
= t^t\eu^{-t}\sqrt t\int_{-\sqrt t}^{\infty}
\exp\Bigl(t\ln\Bigl(1 + \frac u{\sqrt t}\Bigr) - \sqrt
t\,u\Bigr)\dd u,
\]
desde $ x^t = t^t\exp\bigl(t\ln(1 + u/\sqrt t)\bigr)$ e
$\eu^{-x} = \eu^{-t}\eu^{-\sqrt tu}$.

\textbf{2.} Para fixo $ u $ e $ t \to \infty $: $ t\ln(1 +
u/\sqrt t) - \sqrt tu = t\bigl(\frac u{\sqrt t} -
\frac{u^2}{2t} + o(\frac1t)\bigr) - \sqrt tu = -\frac{u^2}2 +
o(1)$: $ g_t(u) \to \eu^{-u^2/2}$.

\textbf{3.} Definir $\psi_1(h) = \varphi(h) + \frac{h^2}4$ sobre
$\intoc{-1}1$: $\psi_1(0) = 0$ e $\psi_1'(h) = \frac1{1+h} -
1 + \frac h2 = \frac{h(h-1)}{2(1+h)}$, que é $\geq 0$ sobre
$\intoc{-1}0$ e $\leq 0$ sobre $\intcc01$: $\psi_1 \leq 0$,
ou seja\ $\varphi(h) \leq -h^2/4$ lá. Definir $\psi_2(h) =
\varphi(h) + ch $ sobre $\intco1\infty $, $ c = 1 - \ln2$: $\psi_2(1)
= \ln2 - 1 + c = 0$ e $\psi_2'(h) = c - \frac h{1+h} \leq c -
\frac12 < 0$: $\varphi(h) \leq -ch $ para $ h \geq 1$. Agora para $ t
\geq 1$: se $\abs u \leq \sqrt t $, $ g_t(u) =
\eu^{t\varphi(u/\sqrt t)} \leq \eu^{-t(u/\sqrt t)^2/4} =
\eu^{-u^2/4}$; se $ u \geq \sqrt t $, então $ t\,\varphi(u/\sqrt t) \leq
-ct\cdot\frac u{\sqrt t} = -c\sqrt t\,u \leq -cu $ (como $ t \geq
1$), então $ g_t(u) \leq \eu^{-cu}$. Por isso $ g_t \leq \eu^{-u^2/4} +
\eu^{-cu}\mathbf 1_{u>0}$, \omterm{def:b3:lebesgue:l1}{integrável}, independente de $ t \geq
1$.

\textbf{4.} DCT: $\int_\R g_t(u)\dd u \to
\int_\R\eu^{-u^2/2}\dd u = \sqrt{2\pi}$
(\cref{ex:b3:product:polar} mais a escala
$ u\mapsto\sqrt2\,u $). Com a pergunta 1:
\[
\Gamma(t + 1) \sim \sqrt{2\pi t}\;\Bigl(\frac
t\eu\Bigr)^t,\qquad
n! \sim \sqrt{2\pi n}\,\Bigl(\frac n\eu\Bigr)^n .
\]

\textbf{5.} De \cref{exo:b3:product:8}, $ W_{2n} =
\frac\pi2\,\frac{(2n)!}{4^n(n!)^2} =
\frac\pi2\,4^{-n}\binom{2n}n $. Stirling:
\[
\binom{2n}{n} = \frac{(2n)!}{(n!)^2}
\sim \frac{\sqrt{4\pi n}\,(2n/\eu)^{2n}}
{2\pi n\,(n/\eu)^{2n}} = \frac{4^n}{\sqrt{\pi n}} .
\]
Então $ W_{2n} \sim \frac12\sqrt{\pi/n}$, consistente com o
recursão $ W_n = \frac{n-1}nW_{n-2}$ (que força $ W_n \sim
W_{n-2}$, e com $ W_nW_{n-1}\cdot n = \frac\pi2$ --- o
relação Wallis clássica --- dá $ W_n \sim
\sqrt{\pi/(2n)}$; os dois assintóticos concordam).

\textbf{6.} $\Gamma(\frac d2 + 1) \sim \sqrt{2\pi\frac
d2}\,(\frac d{2\eu})^{d/2}$, então
\[
v_d = \frac{\pi^{d/2}}{\Gamma(\frac d2 + 1)}
\sim \frac{1}{\sqrt{\pi d}}\Bigl(\frac{2\pi\eu}
d\Bigr)^{d/2} \longrightarrow 0
\]
supergeometricamente (para $ d > 2\pi\eu \approx 17$, cada fator
$< 1$ e diminuindo). Numericamente $ v_1 = 2$, $ v_2 \approx
3.14$, $ v_3 \approx 4.19$, $ v_4 \approx 4.93$, $ v_5 \approx
5.26$, $ v_6 \approx 5.17$: o \omterm{def:b3:rings:primemaximal}{máximo} está em $ d = 5$.

\textbf{7.} Com $ k = \frac n2 + \frac{s\sqrt n}2$ (inteiro,
$ n $ até, $ s $ fixo): pegue logaritmos em $2^{-n}\binom nk =
2^{-n}\frac{n!}{k!(n-k)!}$ e aplique Stirling nos três
fatoriais. Escrita $ k = \frac n2(1 + \varepsilon)$, $ n - k =
\frac n2(1 - \varepsilon)$ com $\varepsilon = s/\sqrt n $:
\[
\ln\Bigl(2^{-n}\binom nk\Bigr)
= -\frac n2\bigl[(1{+}\varepsilon)\ln(1{+}\varepsilon) +
(1{-}\varepsilon)\ln(1{-}\varepsilon)\bigr]
+ \frac12\ln\frac{2}{\pi n(1 - \varepsilon^2)} + o(1),
\]
e o colchete é $\varepsilon^2 + O(\varepsilon^4) =
\frac{s^2}n + O(n^{-2})$: o display tende a $-\frac{s^2}2 +
\frac12\ln\frac2{\pi n}$ até $ o(1)$, ou seja
\[
2^{-n}\binom nk \sim \sqrt{\frac{2}{\pi
n}}\;\eu^{-s^2/2} :
\]
o perfil gaussiano do lançamento de moeda, quantificado ---
de Moivre  a forma local de Laplace, a ser globalizada em
\cref{ch:b3:clt}.

\textbf{8.} (i) DCT converte o limite pontual da questão 2
na convergência das integrais, usando a questão 3
dominador. (ii) A integral gaussiana avalia o limite
$\int\eu^{-u^2/2} = \sqrt{2\pi}$ --- Constante de Stirling
$\sqrt{2\pi}$ \emph{é} a integral gaussiana. (iii) O
substituição $ x = t + \sqrt tu $ é uma mudança afim de
variáveis: invariância de tradução e regra de escala de
Medida Lebesgue (\cref{thm:b3:product:linearchange} em
dimensão $1$).

\textbf{9.} Expanda ambos os termos:
\[
d_n - d_{n+1} = \ln\frac{n!}{(n+1)!} + \Bigl(n +
\frac32\Bigr)\ln(n+1) - \Bigl(n + \frac12\Bigr)\ln n - 1
= \Bigl(n + \frac12\Bigr)\ln\frac{n+1}n - 1,
\]
os termos $-\ln(n+1)$ e $(n + \frac32)\ln(n+1)$ combinando
em $(n + \frac12)\ln(n+1)$.

\textbf{10.} Para $ t = \frac1{2n+1}$: $\frac{1+t}{1-t} =
\frac{2n+2}{2n} = \frac{n+1}n $, e $ n + \frac12 =
\frac1{2t}$; a série estranha $\ln\frac{1+t}{1-t} =
2\sum_{k\geq0}\frac{t^{2k+1}}{2k+1}$ dá
\[
\Bigl(n + \frac12\Bigr)\ln\frac{n+1}n
= \sum_{k\geq0}\frac{t^{2k}}{2k+1}
= 1 + \frac{t^2}3 + \frac{t^4}5 + \cdots
\]
Subtrair $1$. Limite inferior: apenas o primeiro termo,
$\frac{t^2}3 = \frac1{3(2n+1)^2}$. Limite superior: diminuir tudo
denominadores para $3$ e somar as séries geométricas:
$\frac{t^2}{3(1 - t^2)} = \frac1{3((2n+1)^2 - 1)} =
\frac1{12n(n+1)} = \frac1{12n} - \frac1{12(n+1)}$.

\textbf{11.} Somando o limite superior de $ n $ a $\infty $
(com $ d_m \to 0$): $ d_n < \frac1{12n}$. Para o limite inferior:
$\frac1{12m+1} - \frac1{12(m+1)+1} =
\frac{12}{(12m+1)(12m+13)}$, e
\begin{align*}
\frac1{3(2m+1)^2} > \frac{12}{(12m+1)(12m+13)}
&\iff (12m+1)(12m+13) > 36(2m+1)^2 \\
&\iff 168m + 13 > 144m + 36,
\end{align*}
verdadeiro para $ m \geq 1$. Resumindo esta minorante telescópica:
$ d_n > \frac1{12n+1}$. Exponenciar dá o clássico
colchetes de $ n!$.

\textbf{12.} (a) Erro relativo $= \eu^{d_n} - 1 <
\eu^{1/(12n)} - 1 < \frac{1.1}{12n}$ para $ n $ grande; $<
10^{-6}$ assim que $12n \geq 1.1\cdot10^6$, e o declarado
$ n \geq 83\,334$ é suficiente ($\frac1{12n} \leq 10^{-6}$
já implica isso). (b) $\log_{10}(100!) =
\frac12\log_{10}(200\pi) + 200 - 100\log_{10}\eu +
d_{100}\log_{10}\eu = 1.39906 + 200 - 43.42945 + 0.00036
\approx 157.96997$, então $100! \approx 10^{0.96997} \cdot
10^{157} = 9.333\cdot10^{157}$; a janela garantida
$(\eu^{1/1201}, \eu^{1/1200})$ tem largura abaixo $10^{-6}$ em
termos relativos --- uma fórmula ``assiptótica'' que é, pelo menos $ n =
100$, um instrumento de precisão.

\textbf{13.} Escrever $\sin^n = \sin^{n-2} - \sin^{n-2}\cos^2$,
e integre o segundo termo por partes ($ u = \cos\theta $,
$\dd v = \sin^{n-2}\cos\theta\,\dd\theta $, $ v =
\frac{\sin^{n-1}}{n-1}$):
\[
\int_0^{\pi/2}\sin^{n-2}\cos^2 =
\Bigl[\cos\theta\,\frac{\sin^{n-1}\theta}{n-1}\Bigr]_0^{\pi/2}
+ \frac1{n-1}\int_0^{\pi/2}\sin^n = \frac{W_n}{n-1} .
\]
Por isso $ W_n = W_{n-2} - \frac{W_n}{n-1}$, ou seja\ $ W_n =
\frac{n-1}nW_{n-2}$. De $ W_0 = \frac\pi2$, $ W_1 = 1$:
\[
W_{2n} = \frac{(2n-1)!!}{(2n)!!}\cdot\frac\pi2
= \frac\pi2\binom{2n}n4^{-n},
\qquad
W_{2n+1} = \frac{(2n)!!}{(2n+1)!!}
= \frac{4^n}{(2n+1)\binom{2n}n},
\]
convertendo fatoriais duplos por $(2n)!! = 2^nn!$ e
$(2n-1)!! = \frac{(2n)!}{2^nn!}$. Finalmente $ nW_nW_{n-1} =
(n-1)W_{n-1}W_{n-2}$ pela recursão: constante, igual a
$1\cdot W_1W_0 = \frac\pi2$.

\textbf{14.} $ W_{2n+1} \leq W_{2n} \leq W_{2n-1}$ (pontualmente
monotonicidade de $\sin^n $) e $\frac{W_{2n-1}}{W_{2n+1}} =
\frac{2n+1}{2n} \to 1$ espremer $\frac{W_{2n}}{W_{2n+1}} \to
1$. Combinado com $ W_{2n}W_{2n+1} = \frac{\pi}{2(2n+1)}$
(questão 13): $ W_{2n}^2 \sim \frac\pi{4n}$, então $ W_{2n} \sim
\frac12\sqrt{\frac\pi n}$ e $\binom{2n}n4^{-n} =
\frac2\pi W_{2n} \sim \frac1{\sqrt{\pi n}}$.

\textbf{15.} As questões 9 a 10 nunca usaram o valor do
constante: com $ e_n = \ln n! - (n+\frac12)\ln n + n $, o
diferenças $ e_n - e_{n+1}$ deitar em $\bigl(0, \frac1{12n} -
\frac1{12(n+1)}\bigr)$, então $(e_n)$ diminui enquanto $(e_n -
\frac1{12n})$ aumenta: sequências adjacentes, convergindo para um
comum $\ell $. Por isso $ n! \sim K n^{n+1/2}\eu^{-n}$, $ K =
\eu^\ell $.

\textbf{16.} Substituindo a constante desconhecida Stirling em
o binômio central:
\[
\binom{2n}n4^{-n}\sqrt n \sim
\frac{K\,(2n)^{2n+1/2}\eu^{-2n}}{\bigl(K\,n^{n+1/2}
\eu^{-n}\bigr)^2}\,4^{-n}\sqrt n
= \frac{2^{2n}\sqrt{2}\,K\,n^{2n+1/2}}{K^2\,n^{2n+1}}
\,4^{-n}\sqrt n = \frac{\sqrt2}{K},
\]
e pergunta 14 forças $\frac{\sqrt2}K = \frac1{\sqrt\pi}$:
$ K = \sqrt{2\pi}$. As Partes III-IV reprovam assim Stirling de
arranhar; inserindo-o nas avaliações de identidade da Parte I
$\int_\R\eu^{-u^2/2}\dd u = \sqrt{2\pi}$ sem polar
coordenadas: Wallis e Gauss apoiam-se mutuamente.

\textbf{17.} $\frac{\binom{2n}{n+j}}{\binom{2n}n} =
\frac{(n!)^2}{(n+j)!\,(n-j)!} =
\prod_{i=1}^{j}\frac{n-i+1}{n+i}$ para $ j \geq 0$ (e por
simetria para $ j < 0$). Tomando logaritmos, com $1 \leq i
\leq j \leq K\sqrt n $:
\[
\ln\frac{n-i+1}{n+i} = \ln\Bigl(1 - \frac{i-1}n\Bigr) -
\ln\Bigl(1 + \frac in\Bigr) = -\frac{2i-1}{n} +
O\Bigl(\frac{i^2}{n^2}\Bigr),
\]
e $\sum_{i\leq j}(2i - 1) = j^2$, enquanto o erro soma
$ O(j^3/n^2) = O(n^{-1/2})$: uniformemente,
$\exp\bigl(-\frac{j^2}n + O(n^{-1/2})\bigr)$.

\textbf{18.} $\eu^{-n}\frac{n^n}{n!} \sim \eu^{-n}
\frac{n^n}{\sqrt{2\pi n}\,n^n\eu^{-n}} =
\frac1{\sqrt{2\pi n}}$. Uma variável de Poisson de média $ n $ tem
desvio padrão $\sqrt n $, e $\frac1{\sqrt{2\pi n}}$ é
exatamente a altura do pico gaussiano $\frac1{\sigma\sqrt{2\pi}}$:
o CLT local, visualizado no modo.

\textbf{19.} Para $ a \in \intoo01$, o lema da inclinação de
\cref{pb:b3:lebesgue:1} (questão 14 aí), aplicada ao
convexo $\log\Gamma $ em volta $ n $, dá $(n-1)^a \leq
\frac{\Gamma(n+a)}{\Gamma(n)} \leq n^a $: a proporção para $ n^a $
é espremido por $(1 - \frac1n)^a \to 1$. Para $ a = m + a'$
($ m \in \N $, $ a' \in \intco01$): $\Gamma(n+a) = (n + a - 1)
\cdots(n + a')\Gamma(n + a')$, e cada um dos $ m $ fatores é
$ n(1 + O(\frac1n))$: multiplique as estimativas.

\textbf{20.} A recursão $ v_d = \frac{2\pi}dv_{d-2}$
(de $\Gamma(\frac d2 + 1) = \frac d2\Gamma(\frac d2)$)
dá, de $ v_1 = 2$, $ v_2 = \pi $:
\[
v_3 = \frac{4\pi}3,\quad v_4 = \frac{\pi^2}2,\quad
v_5 = \frac{8\pi^2}{15},\quad v_6 = \frac{\pi^3}6,\quad
v_7 = \frac{16\pi^3}{105}.
\]
A proporção $\frac{2\pi}d $ excede $1$ exatamente para $ d \leq 6$,
então cada paridade aumenta e depois diminui; numericamente $ v_4
\approx 4.93$, $ v_5 \approx 5.26$, $ v_6 \approx 5.17$: o
\omterm{def:b3:rings:primemaximal}{máximo} geral é $ d = 5$. Função de geração: $ v_{2k} =
\frac{\pi^k}{k!}$, então $\sum_kv_{2k}x^{2k} =
\eu^{\pi x^2}$ --- todos os volumes de bolas de dimensão uniforme rolados
em um exponencial e um decaimento supergeométrico instantâneo
estimativa para $ v_d $.

\textbf{21.} Stirling no numerador e no denominador, com $ k =
\alpha n $:
\[
\binom n{\alpha n} \sim
\frac{\sqrt{2\pi n}\,n^n}
{\sqrt{2\pi\alpha n}\,(\alpha n)^{\alpha n}\,
\sqrt{2\pi(1-\alpha)n}\,((1-\alpha)n)^{(1-\alpha)n}}
= \frac{\eu^{nH(\alpha)}}{\sqrt{2\pi\alpha(1-\alpha)n}},
\]
desde $ n^n/(\alpha n)^{\alpha n}((1-\alpha)n)^{(1-\alpha)n} =
\alpha^{-\alpha n}(1-\alpha)^{-(1-\alpha)n} =
\eu^{nH(\alpha)}$ (os poderes de $ n $ cancelar: $\alpha n +
(1-\alpha)n = n $) e o $\eu^{-n}$ é cancelado da mesma forma. No
$\alpha = \frac12$: $ H = \ln2$ e o pré-fator é
$\sqrt{2/(\pi n)}$ --- pergunta 5 novamente. Concavidade estrita de
$ H $ (sua segunda derivada $-\frac1{\alpha(1-\alpha)} < 0$)
coloca seu \omterm{def:b3:rings:primemaximal}{máximo} $\ln 2$ no $\alpha = \frac12$ apenas: para
$\alpha \neq \frac12$, $\binom n{\alpha n}2^{-n} \approx
\eu^{-n(\ln2 - H(\alpha))}$ decai exponencialmente --- o
mecanismo combinatório por trás de cada asserção de concentração
sobre lançamentos de moedas.

\textbf{22.} De $ s_{d-1} = dv_d $ e questão 20:
\[
s_0 = 2,\ \ s_1 = 2\pi,\ \ s_2 = 4\pi,\ \ s_3 = 2\pi^2,\ \
s_4 = \frac{8\pi^2}3,\ \ s_5 = \pi^3,\ \ s_6 =
\frac{16\pi^3}{15},
\]
numericamente $2,\ 6.28,\ 12.57,\ 19.74,\ 26.32,\ 31.01,\
33.07$; e $ s_7 = \frac{\pi^4}3 \approx 32.47 < s_6$: o
\omterm{def:b3:rings:primemaximal}{máximo} é o $6$-esfera. A recursão $ s_{d+1} =
(d+2)\,v_{d+2} = (d+2)\,\frac{2\pi}{d+2}\,v_d = 2\pi v_d =
\frac{2\pi}d\,s_{d-1}$ mostra o mesmo $\frac{2\pi}d $-dirigido
ascensão e queda supergeométrica quanto aos volumes: dimensão passada
sete, as esferas encolhem mais rápido do que qualquer geometria
sequência.

\textbf{23.} A pergunta 11 diz exatamente $\frac1{12n+1} < d_n <
\frac1{12n}$, e
\[
\frac1{12n} - \frac1{12n+1} = \frac1{12n(12n+1)} =
O\Bigl(\frac1{n^2}\Bigr),
\]
então $ d_n = \frac1{12n} + O(\frac1{n^2})$ e $\eu^{d_n} = 1 +
\frac1{12n} + O(\frac1{n^2})$; multiplicando por
$\sqrt{2\pi n}(n/\eu)^n $ fornece a fórmula corrigida. No $ n =
10$: $\sqrt{20\pi}\,(10/\eu)^{10} = 7.92665 \times 453999.3
\approx 3\,598\,696$, baixo por $30\,104$ (erro relativo
$8.3\cdot10^{-3}$); multiplicando por $1 + \frac1{120}$ dá
$3\,628\,685$, baixo por $115$ (erro relativo
$3.2\cdot10^{-5}$). O próprio suporte fixa os pinos $10!$ entre
$3\,598\,696\,\eu^{1/121} \approx 3\,628\,559$ e
$3\,598\,696\,\eu^{1/120} \approx 3\,628\,808$ --- a parte superior
limite está errado em oito unidades em sete dígitos.

\textbf{24.} Substituição da Parte I $ x = t + \sqrt t\,u $,
aplicado à integral truncada, dá
\[
\int_0^{t} x^{t}\eu^{-x}\,\dd x
= t^{t}\eu^{-t}\sqrt t\int_{-\sqrt t}^{0} g_t(u)\,\dd u ,
\]
o intervalo $0 \leq x \leq t $ tornando-se $-\sqrt t \leq u \leq
0$. O dominador da questão 3 cobre $ g_t\mathbf 1_{u < 0}$
também, rendimentos de convergência tão dominados
\[
\int_{-\sqrt t}^{0}g_t(u)\,\dd u \longrightarrow
\int_{-\infty}^{0}\eu^{-u^2/2}\dd u = \frac{\sqrt{2\pi}}2 ,
\]
enquanto a pergunta 4 dá $\Gamma(t+1) \sim
t^t\eu^{-t}\sqrt t\,\sqrt{2\pi}$. A proporção tende a
$\frac12$. Probabilisticamente: uma variável aleatória gama de
forma grande coloca assintoticamente metade de sua massa em cada lado
de seu modo --- a simetria do teorema do limite central, leia
fora de uma substituição.

\textbf{25.} Deixar $\lambda = \frac{\alpha}{1-\alpha} \in
\intoc01$. Para $ k \leq \lfloor\alpha n\rfloor \leq \alpha n $
nós temos $\lambda^{k - \alpha n} \geq 1$, por isso
\[
\sum_{k=0}^{\lfloor\alpha n\rfloor}\binom nk
\leq \lambda^{-\alpha n}\sum_{k=0}^{n}\binom nk\lambda^{k}
= \Bigl(\lambda^{-\alpha}(1 + \lambda)\Bigr)^{n},
\]
e com $\lambda = \frac\alpha{1-\alpha}$:
\[
\lambda^{-\alpha}(1+\lambda)
= \alpha^{-\alpha}(1-\alpha)^{\alpha}\cdot\frac1{1-\alpha}
= \alpha^{-\alpha}(1-\alpha)^{-(1-\alpha)}
= \eu^{H(\alpha)} .
\]
Otimalidade: minimizando $ f(\lambda) = -\alpha\ln\lambda +
\ln(1+\lambda)$ sobre $\lambda > 0$, a equação $ f'(\lambda)
= -\frac\alpha\lambda + \frac1{1+\lambda} = 0$ tem o único
solução $\lambda = \frac\alpha{1-\alpha}$, um mínimo desde
$ f'' > 0$ --- a escolha da inclinação exponencial (Chernoff).
Reconciliação: pela questão 21 o termo único $ k =
\lfloor\alpha n\rfloor $ já está em ordem
$\eu^{nH(\alpha)}/\sqrt{2\pi\alpha(1-\alpha)n}$, então
\[
\frac{\eu^{nH(\alpha)}}{C\sqrt n} \leq
\sum_{k\leq\alpha n}\binom nk \leq \eu^{nH(\alpha)} :
\]
a taxa $ H(\alpha)$ é exato, a soma total custando no \omterm{def:b3:rings:primemaximal}{máximo}
um fator $\sqrt n $ ao longo do seu maior prazo. Dividido por $2^n $,
este é o limite da cauda da moeda justa $\P(S_n \leq \alpha n) \leq
\eu^{-n(\ln2 - H(\alpha))}$ --- concentração de medida em
uma linha.
\end{solution}
